\section{Constructing the inert comparison in every parity}
\label{sec:sr-inert-matching-construction}

We construct the finite matching maps at an inert place outside
the prescribed ramification set. The ranks of the two Hermitian
spaces need not be even. This distinction matters: results stated
only for even Hermitian rank do not supply this clause of the packet.

Write \(k=F_v\), \(L=E_w\), and let \(\overline{\phantom{x}}\)
denote the nontrivial automorphism of the unramified quadratic
extension \(L/k\). Put
\[
 d=a+a'=2n,\qquad a>a'\ge1,
 \qquad K_H=U(\Lambda_H),\quad K=U(\Lambda_W),
\]
where \(H\) is either \(V\) or \(V'\), and \(\Lambda_H,\Lambda_W\)
are the actual localizations of the prescribed lattices.
The additive character is the prescribed \(\psi_v\).
Haar measure on each compact group in this section has total mass one.

\subsection{Why the actual local data are unramified}
The base hypotheses at an inert place give odd residue characteristic,
residue cardinality greater than three, and
\[
 \{x\in k:\psi_v(x\mathcal O_k)=1\}=\mathcal O_k.
\]
For the standard global character the group on the left is
\(\mathfrak D_{k/\mathbb Q_p}^{-1}\). This follows directly by
localizing the trace-dual definition of the codifferent, or by
writing the character as the standard character of \(\mathbb Q_p\)
composed with the trace. Consequently the different exponent of
\(F/\mathbb Q\) at \(v\) is zero. The tower formula for differents
and the unramifiedness of \(E/F\) give
\(\mathfrak D_{L/\mathbb Q_p}=\mathcal O_L\).
This checks a potentially relevant normalization: the source's
skew lattice is initially defined using the absolute trace dual.
Here that dual is the ordinary \(\mathcal O_L\)-dual.
It is therefore a genuinely unimodular skew-Hermitian lattice,
just as the two Hermitian lattices are unimodular.

Choose a trace-zero unit \(\delta\in\mathcal O_L^\times\).
Multiplying a skew-Hermitian form by \(\delta^{-1}\) makes it
Hermitian and changes none of these duality assertions.
Such a unit exists by lifting an element of the residue field
whose conjugate differs from itself; subtract its conjugate.
The norm map on units is surjective. One way to verify this is
to use the surjectivity of the norm between the two finite
residue fields and then lift successively: the first-order norm
of \(1+\varpi^r x\) is \(1+\varpi^r\operatorname{Tr}(x)\).
The residue trace is onto, so each correction can be made.

Reduction modulo \(\varpi\), followed by Hensel lifting, now
splits a hyperbolic plane off every even unimodular Hermitian
lattice. To give the lifting step explicitly, start with a primitive
isotropic residue vector, lift it, choose a partner whose pairing
is a unit, and adjust along that partner to make its norm zero.
The equation has a unit derivative because two is a unit.
Normalize the pairing and replace the partner by an isotropic
partner by subtracting half its norm times the first vector.
The orthogonal complement is again unimodular.
Induction gives a hyperbolic basis of \(\Lambda_W\), since its
rank \(2n\) is even. In particular \(U(W)\) is the split-rank
unitary group \(U(n,n)(k)\), with its specified hyperspecial subgroup
\(K\). For either parity of \(a\) and \(a'\), the same argument
gives the usual unimodular local Hermitian models of \(V,V'\).

Let \(\alpha_v,\beta_v\) be the two local splitting characters.
They are unramified and their restrictions to \(k^\times\)
are respectively \(\eta^{a'}\) and \(\eta^a\), where \(\eta\)
is the unramified quadratic character of \(k^\times\).
The common uniformizer of \(k\) and \(L\) therefore satisfies
\[
 \alpha_v(\varpi)=(-1)^{a'}=(-1)^a=\beta_v(\varpi).
\]
Their values on units are one, so
\begin{equation}\label{eq:sr-inert-character-equality}
 \alpha_v=\beta_v=\chi,\qquad \chi^2=1,\qquad
                 \alpha_v/\beta_v=1.
\end{equation}
Thus the first member of each dual pair has trivial auxiliary
character, and the second member has the same unramified splitting
character \(\chi\). This verifies the splitting data, including
the case of odd \(a,a'\).

\subsection{Transport to the ordinary local model}
Use a hyperbolic polarization \(W=Y\oplus Y^\vee\) with
\(\Lambda_W=\Lambda_Y\oplus\Lambda_{Y^\vee}\).
For \(H=V,V'\), put
\[
 \mathcal S_H=\mathcal S(H\otimes_LY^\vee)
              \simeq\mathcal S(H^n),\qquad
                     e_H=1_{\Lambda_H^n}.
\]
All additive measures are self-dual for \(\psi_v\).
Because the tensor lattice is selfdual, its volume is one
in dual coordinates; the Fourier transform carries \(e_H\)
to the corresponding dual indicator.
Integral unipotents fix \(e_H\), since their quadratic phases
are integral and two is a unit. Integral Levi matrices fix it
by the determinant formula; their splitting-character factor is one.
The Fourier generators fix it as well, by the preceding selfdual
Fourier identity and the conductor-zero, odd-residue normalization.
These are the lattice-model generators of both hyperspecial groups,
so \(e_H\) is fixed by \(K_H\times K\).

This model is the actual model in different coordinates, not an
additional choice of an unrelated Weil representation.
Indeed a lattice frame of \(\Lambda_H,\Lambda_W\) identifies their
tensor lattice with integral matrices and identifies the symplectic
pairing with the trace pairing just used. The Stone--von Neumann
intertwiner transports the prescribed Heisenberg operators to
ordinary Schr\"odinger translations and modulations.
The phase in the prescribed vacuum condition is one here:
the tensor-lattice pairings and the integral element \(\theta\)
occurring in that condition are integral, and two is a unit.
Its lattice-fixed line is one-dimensional: modulation by the
integral lattice forces support in its annihilator, and integral
translations then force a constant on that lattice.
Rescale the intertwiner to send the prescribed nonzero vacuum
exactly to \(e_H\). This also handles any original normalization
of that vacuum.

Egorov covariance implies that the transported Weil operators
differ from the ordinary ones by a character of
\(U(H)\times U(W)\). The first member preserves the whole
Schr\"odinger polarization \(H\otimes_LY^\vee\).
The specified Siegel rule on this polarization fixes its action
for every \(g\in U(H)\): it is pullback, since its determinant
has norm one and its auxiliary character is one.
On the second member the same rule fixes the scalar discrepancy
on the Siegel parabolic \(P\), including its Levi.
A character trivial on \(P\) is trivial on its opposite as well,
because characters are conjugation invariant.
The two opposite parabolics generate \(U(W)\):
their root groups generate the special unitary subgroup, and
the determinant subgroup is already represented by the Levi
(the latter assertion is Hilbert's theorem 90).
Thus the scalar discrepancy is one. This argument includes
\(U(1,1)\); for a compact rank-one first member its entire action
was already fixed by the polarization rule.
The source and ordinary actions, their compact stabilizers
and their vacua have all been identified simultaneously.

Evaluation at zero now supplies a nonzero first-member-invariant
linear functional
\[
       e_H^0(\phi)=\phi(0),\qquad e_H^0(e_H)=1.
\]
The invariance uses \eqref{eq:sr-inert-character-equality}:
the first member acts by pullback with no remaining character.
Let \(P\) be the Siegel parabolic of \(U(W)\), stabilizing \(Y\).
For a Levi element with action \(t\in\mathrm{GL}_n(L)\) on \(Y\),
the prescribed Siegel rule gives
\begin{equation}\label{eq:sr-inert-parabolic-functional}
 e_H^0\bigl(\omega_H(p)\phi\bigr)
   =\chi(\det t)|\det t|_L^{\dim_L H/2}\,e_H^0(\phi).
\end{equation}
The unipotent radical has scalar one in this formula.

\subsection{The two complementary spherical functions coincide}
Set \(r_a(p)=\chi(\det t)|\det t|_L^{a/2}\) and
\(r_{a'}(p)=\chi(\det t)|\det t|_L^{a'/2}\).
The modular character of \(P\) is
\[
                   \delta_P(p)=|\det t|_L^n.
\]
For example the unipotent radical is the additive space of
Hermitian \(n\)-by-\(n\) matrices; conjugation by \(t\)
has \(k\)-determinant of absolute value
\(|\det t|_L^n\), as is seen first for diagonal \(t\)
and then by multiplicativity. Hence
\begin{equation}\label{eq:sr-inert-complementary-character}
                          r_a r_{a'}=\delta_P.
\end{equation}

Define
\[
 f_a(h)=e_V^0(\omega_V(h)e_V),\qquad
 f_{a'}(h)=e_{V'}^0(\omega_{V'}(h)e_{V'}).
\]
They are right \(K\)-invariant, equal to one on \(K\), and
satisfy \(f_a(ph)=r_a(p)f_a(h)\), with the analogous rule
for \(a'\). The Iwasawa decomposition \(U(W)=PK\)
therefore determines each function uniquely.
Let \(I_a,I_{a'}\) be the smooth induced spaces with these
transformation rules. The compact-picture bilinear pairing
\[
                    B(F,F')=\int_K F(k)F'(k)\,dk
\]
is invariant under simultaneous right translation.
To check the assertion, decompose \(kh=p(k,h)k_h\).
The change of variable on \((P\cap K)\backslash K\)
has density \(\delta_P(p(k,h))^{-1}\).
The product of the two transformation factors is
\(\delta_P(p(k,h))\) by
\eqref{eq:sr-inert-complementary-character}, and cancels it.
The inner integration over \(P\cap K\) has probability mass one.
This proves the assertion directly, also at reducible induction
parameters.

Normalize the compact averages of the evaluation functionals by
\[
 \ell_a(\phi)=\int_K e_V^0(\omega_V(k)\phi)\,dk,\qquad
 \ell_{a'}(\phi')=\int_K e_{V'}^0(\omega_{V'}(k)\phi')\,dk.
\]
Each integral is a finite average for every smooth vector.
Both functionals annihilate the first-member augmentation,
are \(K\)-invariant, and take the value one on their vacua.
Their spherical coefficients are
\[
 \Omega_a(h)=\ell_a(\omega_V(h)e_V)
              =B(R(h)f_a,f_{a'}),\qquad
 \Omega_{a'}(h)=B(f_a,R(h)f_{a'}).
\]
The invariant pairing proves
\(\Omega_a(h)=\Omega_{a'}(h^{-1})\).

For completeness the last inverse may be removed using the
specified hyperspecial group, without changing that group.
Integral hyperbolic elimination gives
\[
 h=k_1\operatorname{diag}
   (\varpi^{\lambda_1},\ldots,\varpi^{\lambda_n},
    \varpi^{-\lambda_1},\ldots,\varpi^{-\lambda_n})k_2,
 \quad k_1,k_2\in K.
\]
One obtains this by diagonalizing the relative position of
the two selfdual lattices \(\Lambda_W,h\Lambda_W\):
their elementary divisors occur in opposite pairs.
At each step a primitive extremal vector and a dual partner
are carried to a hyperbolic pair by integral isometries,
and the orthogonal complement gives the next step.
The integral Weyl element exchanging the two vectors of
every hyperbolic pair conjugates the displayed diagonal
to its inverse. Thus \(h^{-1}\in KhK\).
The two coefficients are bi-\(K\)-invariant, and consequently
\begin{equation}\label{eq:sr-inert-common-spherical-function}
                    \Omega_a(h)=\Omega_{a'}(h)=:\Omega(h),
                       \qquad\Omega(1)=1.
\end{equation}

\subsection{Constructing an actual irreducible common image}
We use two precise published local facts. For unramified quadratic
unitary dual pairs with selfdual lattices in odd residue
characteristic, the joint hyperspecial fixed space is generated
from the lattice vacuum by either spherical Hecke algebra;
an irreducible spherical quotient receives a nonzero vacuum
and has a one-dimensional hyperspecial fixed line.
These are Howe's theorems as stated in Harris,
\emph{Cohomological automorphic forms on unitary groups, II},
Appendix~2.A, Theorems~2.A.1--2.A.2.
\footnote{\url{https://webusers.imj-prg.fr/~michael.harris/periodrelations.pdf}}
For a nonzero big theta lift of an irreducible admissible
representation, finite length and Howe duality give a nonzero
irreducible quotient. We use Gan--Takeda,
\emph{A proof of the Howe duality conjecture},
Theorem~1.2 and its introductory finite-length assertion.
\footnote{\url{https://arxiv.org/abs/1407.1995}}
Their field, lattice, residue, character and splitting hypotheses
were checked above; no even-rank restriction is being added.

Here is the exact quotient argument needed from those facts.
Let
\[
 Q_a=\mathcal S_V/
    \langle\omega_V(g,1)\phi-\phi:g\in U(V),\ \phi\in\mathcal S_V\rangle .
\]
Evaluation at zero proves \(Q_a\ne0\).
Moreover
\begin{equation}\label{eq:sr-inert-coinvariant-fixed-line}
                            Q_a^K=\mathbb C[e_V].
\end{equation}
Indeed compact averaging is exact on smooth complex
representations. A \(K\)-fixed vector of \(Q_a\) therefore
has a \(K\)-fixed lift; average this lift over \(K_V\) as well.
The second averaging does not change its image, because the
whole first member acts trivially on \(Q_a\).
The resulting joint fixed vector is a Hecke translate of
\(e_V\) under the first spherical Hecke algebra.
On first-member coinvariants a Hecke function acts by
its integral, so its image is a scalar multiple of \([e_V]\).
This proves \eqref{eq:sr-inert-coinvariant-fixed-line}.

Let \(q_a:Q_a\twoheadrightarrow\pi_a\) be the nonzero
irreducible theta quotient. Its vacuum is nonzero by the cited
unramified theorem. Equation
\eqref{eq:sr-inert-coinvariant-fixed-line} and exactness of
\(K\)-averaging give \((\ker q_a)^K=0\).
The \(K\)-invariant functional \(\ell_a\) annihilates this
kernel: on a kernel vector its value is unchanged by
\(K\)-averaging, and that average is zero.
Thus \(\ell_a\) descends to a nonzero functional on \(\pi_a\).
The coefficient map
\[
 C_a(\phi)(h)=\ell_a(\omega_V(h)\phi)
\]
therefore factors through \(\pi_a\). Its induced map is
nonzero, since \(C_a(e_V)(1)=1\), so is injective by
irreducibility. Its image is an irreducible representation
and is generated by \(C_a(e_V)=\Omega\), since every nonzero
vector in an irreducible representation is cyclic.
The same reasoning applies to \(a'\).

It follows from \eqref{eq:sr-inert-common-spherical-function}
that both coefficient maps have the same actual image
\[
            \Pi_v=
              \operatorname{span}_{\mathbb C}
                 \{h\longmapsto\Omega(hh_0):h_0\in U(W)\}.
\]
Set \(p_v=C_a\) and \(p'_v=C_{a'}\), with codomain \(\Pi_v\).
They are surjective and annihilate their respective
first-member augmentations. For every \(h_0\in U(W)\),
\[
 C_a(\omega_V(h_0)\phi)(h)=C_a(\phi)(hh_0),
\]
and the same equality holds for \(C_{a'}\).
Thus they induce the same second-member action, not merely
isomorphic abstract spaces. Their vacuum images are the
same nonzero function \(\Omega\). Finally every equivariant
endomorphism of \(\Pi_v\) is scalar: it preserves the
one-dimensional \(K\)-fixed line, so acts by a scalar on
\(\Omega\), and hence on all its translates.

Transport \(p_v,p'_v\) back by the vacuum-normalized local
intertwiners above. This constructs exactly the inert
finite matching datum for the prescribed models, including
surjectivity, augmentation, common action, irreducibility,
scalar endomorphisms, and equality of the actual vacua.

The positive-rank statement covers every application here:
the packet inequalities give \(a'\ge2\) and \(n\ge1\).
If a zero Hermitian rank is included in an auxiliary formulation,
its Schwartz space is \(\mathbb C\), its Weil action and
normalized vacuum coefficient are trivial, and its common
image is the constant functions. The same compact pairing
with \(r_0=1\) gives the constant coefficient on its
complementary positive-rank side. If all ranks are zero
both maps are the identity of \(\mathbb C\).
No \(\mathrm{GL}_0\) stabilizer or positive-rank Howe assertion
is used for those trivial spaces.

\subsection{Scope and simultaneous choices}
Together with the split-unramified and ramified constructions,
this proves the finite matching clause at every finite place
of a fixed actual triple. At a place outside the ramification
set, normalize by the prescribed vacuum, as above.
At a ramified place use the fixed full-rank local normalization
already constructed. The resulting local maps may all be
chosen at once; each codomain is the specified common
function image, so no incompatible choice of quotient
isomorphisms is needed. Their tensor normalization is fixed
before choosing polynomial vectors or integral theta tests.
This assertion does not impose a uniform norm estimate on
these maps as the global triple varies.

The archimedean common Fock quotient, the global nonzero
comparison constant and the two codimension-two recursion
identities still have to be supplied to conclude the full
global matching and theta-integral assertions. The construction
above closes the inert finite clause itself; it does not
replace any of those remaining global assertions.
